\section{Literature Review}
\label{sec:literature}

This paper draws on several distinct strands of the finance literature.
We synthesize insights from structural credit models, closed-end fund pricing, convertible bond theory, cryptocurrency finance, and capital structure research to build a framework suited to the unique characteristics of Strategy Inc (formerly MicroStrategy) as a BTC-backed equity vehicle.

\subsection{Structural Credit Models}

The foundational framework for linking firm asset values to security prices originates with \citet{merton1974pricing}, who models equity as a European call option on the firm's assets with a strike equal to the face value of debt.
This insight---that limited liability creates option-like payoffs for equity holders---is central to our model's treatment of NAV with a limited liability floor: $\text{NAV}_t = (H_t S_t - D_t)^+$.

\citet{black1976valuing} extend Merton's framework to incorporate safety covenants and subordination, both relevant to Strategy's multi-layered capital structure with preferred stock tranches of differing seniority.
\citet{leland1994corporate} and \citet{leland1996optimal} develop models with endogenous bankruptcy and optimal capital structure, introducing the trade-off between tax shields and distress costs that informs our survival probability analysis.
More recently, \citet{huang2004structural} provide empirical evidence on the performance of structural models, finding that they substantially underpredict credit spreads---the so-called ``credit spread puzzle''---suggesting that factors beyond pure leverage drive credit risk.

In our setting, the single-asset nature of Strategy's balance sheet (essentially only BTC) makes structural models particularly apt: the firm's asset volatility is directly observable via Bitcoin price dynamics, unlike diversified corporations where asset volatility must be inferred.
However, the presence of multiple layers of preferred stock with distinct priority claims extends the standard two-class (debt/equity) Merton framework into a multi-tranche setting analogous to structured credit.

\subsection{NAV Premium and Discount Dynamics}

Strategy trades at a persistent premium to its Bitcoin net asset value, a phenomenon with close parallels to the closed-end fund (CEF) pricing puzzle.
\citet{lee1991investor} propose that investor sentiment drives CEF discounts and premia, with noise trader risk preventing arbitrage from eliminating mispricings.
\citet{pontiff1995closed} documents that CEF premia predict future returns, consistent with mean-reversion---the empirical pattern we capture through our Ornstein--Uhlenbeck premium process $\pi_t$.

\citet{cherkes2009closed} offer a liquidity-based explanation: CEFs provide access to illiquid assets, justifying a premium when the underlying is hard to access directly.
This theory has particular relevance to Strategy, which offers leveraged, operationally-managed BTC exposure that individual investors may find difficult or costly to replicate, especially pre-2024 when spot Bitcoin ETFs were unavailable in the United States.
\citet{dimson1999closed} provide a comprehensive survey of CEF pricing anomalies, noting that premia tend to be larger for funds holding exotic or hard-to-value assets---a characterization that arguably applies to Bitcoin.

The introduction of spot Bitcoin ETFs in January 2024 \citep{alexander2024bitcoin} represents a structural break in the accessibility of BTC exposure, potentially compressing Strategy's premium over time.
Our mean-reverting premium model with regime-dependent long-run mean $\theta$ can accommodate such shifts.
The premium dynamics we observe---with $\pi_t$ positively correlated with BTC returns ($\rho > 0$)---echo the procyclical sentiment patterns documented in the CEF literature.

\subsection{Bitcoin and Cryptocurrency Finance}

The academic literature on Bitcoin pricing has grown rapidly since \citet{nakamoto2008bitcoin}.
\citet{liu2021risks} establish that cryptocurrency returns are driven by network and momentum factors distinct from traditional asset pricing factors, implying that BTC exposure introduces novel risk dimensions to Strategy's equity.
\citet{makarov2020trading} document arbitrage frictions in cryptocurrency markets that can amplify volatility, relevant to our BTC price dynamics specification.

On volatility modeling, \citet{katsiampa2017volatility} compares GARCH specifications for Bitcoin and finds that models with both short- and long-run volatility components best capture BTC dynamics.
\citet{biais2023equilibrium} provide a general equilibrium model of Bitcoin pricing that links fundamentals (mining costs, network adoption) to price, offering theoretical grounding for the drift parameter $\mu_S$ in our $\mathbb{P}$-measure BTC dynamics.
The jump-diffusion specification we employ for BTC prices under $\mathbb{Q}$ follows \citet{merton1976option} and \citet{kou2002jump}, accommodating the well-documented fat tails and discontinuous price movements in cryptocurrency markets \citep{griffin2020bitcoin}.

\subsection{Bitcoin as a Corporate Treasury Asset}

Strategy's pioneering adoption of Bitcoin as a treasury reserve asset, beginning in August 2020, has spawned a small but growing literature.
\citet{vanelli2024bitcoin} study MicroStrategy's treasury strategy and its effects on firm valuation, documenting the reflexive relationship between BTC acquisitions, equity premium, and subsequent capital raises that we formalize as the ``flywheel'' mechanism.
\citet{amin2024corporate} examine more broadly whether corporate cryptocurrency holdings affect firm value, finding heterogeneous effects depending on the scale of holdings relative to firm size.

Our model's BTC holding dynamics---with accumulation rate $\alpha$ driven by the premium $\pi_t^+$ and discrete financing jumps---formalizes the strategic feedback loop whereby a high premium enables accretive capital raises, which increase BTC per share, which supports the premium.
This is a novel contribution: while practitioner analyses have described this flywheel qualitatively, we provide the first stochastic model of its dynamics suitable for Monte Carlo simulation and quantitative indicator extraction.

\subsection{Convertible Bond Pricing and Hybrid Securities}

Strategy's capital structure includes approximately \$8.2 billion in convertible notes alongside five series of preferred stock with varying terms.
The convertible bond pricing literature, beginning with \citet{brennan1977convertible} and \citet{ingersoll1977contingent}, provides the theoretical basis for valuing these hybrid securities.
\citet{tsiveriotis1998valuing} extend the framework to incorporate credit risk, decomposing convertible value into debt and equity components---an approach we use to define ``effective debt'' $D_t^{\text{eff}}$ that adjusts for deep-in-the-money convertibles.
\citet{ammann2003pricing} provide empirical evidence that convertible bonds are systematically underpriced, a finding with implications for the true leverage embedded in Strategy's capital structure.

The preferred stock layer adds further complexity.
\citet{kallberg2013preferred} document that preferred stocks exhibit hybrid risk characteristics, with returns driven by both interest rate and equity market factors.
Strategy's preferred tranches span from pure fixed-income (STRF, STRD) to equity-linked (STRK with conversion into MSTR common), with the variable-rate STRC representing an innovative instrument that adjusts dividends to stabilize price near par---a mechanism without close precedent in the academic literature.

\subsection{Capital Structure and Leverage}

The classical capital structure theories of \citet{modigliani1958cost} and \citet{myers1984capital} provide the backdrop against which Strategy's aggressive leveraging of a single volatile asset must be understood.
\citet{baker2002market} document that firms time equity issuance to exploit high valuations---precisely the behavior we model through premium-dependent financing intensity $\lambda_M(\pi_t)$.

Strategy's capital structure is unusual in several respects that extend standard theory:
(i)~the firm's assets are a single, publicly-traded, highly volatile cryptocurrency;
(ii)~leverage is deliberately and systematically increased to accumulate more of the same asset;
(iii)~the financing flywheel creates a reflexive feedback loop between equity valuation and asset accumulation;
and (iv)~the multi-layered preferred stock structure creates a priority waterfall with complex interactions between dividend obligations, conversion optionality, and residual equity claims.

These features motivate our multi-tranche extension of the Merton framework with endogenous asset accumulation, a contribution at the intersection of structural credit modeling and corporate treasury management.

\subsection{Simulation Methods}

Our quantitative indicators---including the Implied Funding Requirement Distribution (IFRD), per-share Implied BTC Growth Rate (IBGR), and survival probabilities---require Monte Carlo simulation of correlated stochastic processes.
We follow the methods described in \citet{glasserman2004monte}, employing Cholesky decomposition for correlated Brownian motions and thinning algorithms for state-dependent Poisson jump processes.
The Ornstein--Uhlenbeck process for the log-premium follows the exact discretization of \citet{vasicek1977equilibrium}, ensuring that our discrete-time simulation preserves the known transition density of the continuous-time process.
